\documentclass[../../../include/open-logic-section]{subfiles}

\begin{document}

\olfileid{sth}{story}{blv}

\olsection{అనుబంధం: ఫ్రేగె ప్రాథమిక నియమం V}

\olref[rus]{sec}లో రసెల్ తన
వైరుధ్యాన్ని, ఫ్రేగె తన
\emph{Grundgesetze}లో వివరించిన
వ్యవస్థకు సమస్యగా రూపొందించాడని
చెప్పాం. ఫ్రేగె వ్యవస్థలో
అమాయక ధర్మసంగ్రహానికి
నేరుగా ఇచ్చిన రూపం లేదు.
అందువల్ల ఈ అనుబంధంలో
ఆ వ్యవస్థలో \emph{ఏమి}
ఉందో, అది అమాయక
ధర్మసంగ్రహంతోనూ రసెల్
వైరుధ్యంతోనూ ఎలా సంబంధం
కలిగి ఉందో చాలా క్లుప్తంగా
వివరిస్తాం.

ఫ్రేగె వ్యవస్థ \emph{రెండవ-క్రమానికి}
చెందినది; అది \emph{భావన
విస్తారం} అనే భావాన్ని
రూపొందించడానికి ఉద్దేశించబడింది.\footnote{కచ్చితంగా
చెప్పాలంటే, ఫ్రేగె మరింత సాధారణమైన
భావాన్ని, ప్రమేయపు ``విలువల
వ్యాప్తి''ని, ఔపచారికీకరించడానికి
ప్రయత్నించాడు. భావనల విస్తారాలు
ఆ సాధారణ భావానికి ప్రత్యేక
సందర్భాలు. వివరాలకు
\citet[pp.\ 8--9]{Heck2012} చూడండి.}
ఫ్రేగె నుంచి స్ఫూర్తి పొందిన
సంకేతలిపిని వాడుతూ,
\emph{$F$ భావన విస్తారాన్ని}
$\fregeext{x}{F(x)}$ అని రాస్తాం.
ఈ సాధనం ``$F$'' అనే
\emph{విధేయాన్ని} తీసుకొని,
దానిని ``$\fregeext{x}{F(x)}$''
అనే (ప్రథమ-క్రమ) \emph{పదంగా}
మారుస్తుంది. ఈ సాధనంతో
ఫ్రేగె సభ్యత్వానికి ఈ
\emph{నిర్వచనం} ఇచ్చాడు:
\[
	a \in b =_\text{df} \exists G(b = \fregeext{x}{G(x)} \land Ga)
\]
స్థూలంగా, $a \in b$
అంటే $a$ ఒక భావన కిందికి
వస్తుంది, దాని విస్తారం
$b$ అవుతుంది. (ఇక్కడ
``$\exists G$'' పరిమాణకారకం
రెండవ-క్రమానికి చెందిందని
గమనించండి.) \emph{ప్రాథమిక
నియమం V}గా తెలిసిన ఈ
సూత్రాన్ని కూడా ఫ్రేగె
అంగీకరించాడు:
$$\fregeext{x}{F(x)} = \fregeext{x}{G(x)} \liff \forall x (Fx \liff Gx)$$
స్థూలంగా, రెండు భావనలు
ఒకే వస్తువులకు వర్తించినప్పుడే
వాటి విస్తారాలు ఒకటే.
(ఇక్కడ కూడా ``$F$'',
``$G$'' రెండూ విధేయ
స్థానాల్లో ఉన్నాయి.) ఇప్పుడు
ఒక సరళమైన సూత్రం సభ్యత్వాన్ని
ధర్మ-సంతృప్తితో కలుపుతుంది:

\begin{lem}[\emph{Grundgesetze}లో]\ollabel{lem:Fregeextensions}
$\forall F \forall a(a \in \fregeext{x}{F(x)} \liff Fa)$
\end{lem}

\begin{proof} 
$F$, $a$లను స్థిరపరచండి.
సభ్యత్వ నిర్వచనం ప్రకారం
$a \in \fregeext{x}{F(x)}$
అప్పుడే మరియు అప్పుడే
$\exists G(\fregeext{x}{F(x)}
= \fregeext{x}{G(x)} \land Ga)$;
ప్రాథమిక నియమం V ప్రకారం
అప్పుడే మరియు అప్పుడే
$\exists G(\forall x(Fx \liff Gx) \land Ga)$;
ప్రాథమిక రెండవ-క్రమ
తర్కం ప్రకారం అప్పుడే
మరియు అప్పుడే $Fa$.
\end{proof}

దీని నుంచి అమాయక
ధర్మసంగ్రహం దాదాపు వెంటనే
వస్తుంది:

\begin{lem}[\emph{Grundgesetze}లో.]
$\forall F \exists s \forall a (a \in s \liff Fa)$
\end{lem}

\begin{proof}
$F$ను స్థిరపరచండి;
\olref{lem:Fregeextensions} నుంచి
$\forall a (a \in
\fregeext{x}{F(x)} \liff Fa)$
వస్తుంది; కాబట్టి అస్తిత్వ
సాధారణీకరణ ద్వారా
$\exists s\forall a(a \in s \liff Fa)$.
$F$ యాదృచ్ఛికంగా
ఎంచుకున్నది కాబట్టి
ఫలితం వస్తుంది.
\end{proof}

$\forall x(Fx \liff x \notin x)$
అని నిర్ణయించే $F$ను
తీసుకుంటే రసెల్ వైరుధ్యం
వస్తుంది.

\end{document}
